\subsection{Matched Filter and Detection}\label{sec:detection}

Let $m(\tau)$ be one recorded channel over a leg, expressed in the normalised coordinate~\eqref{eq:tau}. The matched filter correlates it with the two basis functions centred at the lag $t$ (in CPA distances):

\begin{equation}
    y_i(t) = \int m(\tau)\,\Phi_i(\tau - t)\,\mathrm{d}\tau, \quad i = 1,2, \qquad \mathbf{y}(t) = \big(y_1(t),\, y_2(t)\big)^{\top}
    \label{eq:filter_inf}
\end{equation}

In the implementation the integral is a sum over samples $m_k$ taken at $\tau_k = k/(f_s T_s)$, with templates sampled on the same grid and normalised to unit Euclidean norm. For a channel~\eqref{eq:channel} and an unbounded record, \eqref{eq:filter_inf} is linear in $(\cos \beta, \sin \beta)$:

\begin{equation}
    \mathbf{y}(t) = a\,\mathbf{C}(t)^{\top}\begin{pmatrix}\cos\beta \\ \sin\beta\end{pmatrix}, \qquad \mathbf{C}(t) = \begin{pmatrix} c_{11} & c_{12} \\ c_{21} & c_{22} \end{pmatrix}, \qquad c_{ij}(t) = \int \Phi_i(\tau)\,\Phi_j(\tau - t)\,\mathrm{d}\tau.
    \label{eq:filter_out_decomposed}
\end{equation}

$c_{ij}(t)$ is the correlation of $\Phi_i$ with $\Phi_j$ shifted by $t$. The four functions have closed forms (Appendix~\ref{app:correlations}):

\begin{equation}
    c_{11}(t) = c_{22}(t) = \frac{4}{4 + t^2}, \qquad c_{21}(t) = -c_{12}(t) = \frac{2t}{4 + t^2}.
    \label{eq:cij}
\end{equation}

The identity $c_{21} = -c_{12}$ follows from parity alone ($\Phi_1$ even, $\Phi_2$ odd). The identity $c_{22} = c_{11}$ is the consequence of the Hilbert relation announced in Section \ref{subsec:cable_obf}: $\Phi_1$ and $\Phi_2$ have the same autocorrelation. Together they make $\mathbf{C}(t)$ a rotation matrix multiplied by the scalar $2/\sqrt{4 + t^2}$, and a rotation does not change the length of the vector $(\cos\beta, \sin\beta)$. Hence

\begin{equation}
\|\mathbf{y}(t)\|^2 = y_1^2 + y_2^2 = \frac{4a^2}{4 + t^2} \qquad \text{for every phase } \beta.
\label{eq:energy_invariance}
\end{equation}

The energy is maximal at $t = 0$, the CPA, falls to half at $|t| = 2$, and does not depend on which channel is used. With an unbounded record the energy peak therefore locates the crossing for every channel and every geometry.

\paragraph{Detection.} The detector thresholds the energy normalised by the noise variance:

\begin{equation}
\Lambda(t) = \frac{\|\mathbf{y}(t)\|^2}{\sigma_n^2} \;\overset{\mathcal{H}_1}{\underset{\mathcal{H}_0}{\gtrless}}\; \eta.
\label{eq:energy_inf}
\end{equation}

$\mathcal{H}_0$ is the hypothesis ``no cable'' and $\mathcal{H}_1$ ``cable''. $\sigma_n$ is the noise standard deviation per sample and $\eta$ the threshold. Under white Gaussian noise, \eqref{eq:energy_inf} is the GLRT for a signal in $\mathcal{S}$ with unknown amplitudes~\cite{hleis2026obf,kay1998detection,scharf1994matched}. With unit-norm, orthogonal templates, $\mathbf{y}$ is then Gaussian with covariance $\sigma_n^2 \mathbf{I}$ under $\mathcal{H}_0$, so $\Lambda$ follows a $\chi^2$ law with 2 degrees of freedom and the false-alarm probability at one lag is $P_{\mathrm{fa}} = \exp(-\eta/2)$ (Appendix~\ref{app:h0}). In practice the threshold is calibrated empirically instead.